\section{Class-DE Resonant Conversion}
\label{sec:bg-classde}

The isolated path comprises the half bridge, series resonant network, transformer, and secondary rectifier. The rectangular switch-node voltage excites approximately sinusoidal tank current. The rectifier charges the secondary bulk capacitor, whose voltage is regulated by the controller.

For series inductance \(\LRES\) and capacitance \(\CRES\),
\begin{equation}
f_0=\frac{1}{2\pi\sqrt{\LRES\CRES}}.
\label{eq:bg-series-resonance}
\end{equation}
The first-harmonic approximation represents the bridge by its fundamental voltage and the rectifier/load by an equivalent resistance. It supports initial gain and component sizing for approximately sinusoidal tank current. Burst transitions, discontinuous rectifier conduction, and no-load operation violate its steady-state load assumptions.

Discrete inductance and transformer leakage form the series branch; magnetizing inductance remains a separate branch. Although omitted from the first series-resonance estimate, magnetizing current contributes to primary RMS current, loss, and commutation, particularly at light load.

Above resonance, the loaded series tank presents an inductive impedance, providing lagging current for switch-node commutation during dead time. Pump current also contributes at the shared node and must be included in the soft-switching assessment.

The reference LED-driver approximation neglects magnetizing current, whereas an LLC gain model retains finite magnetizing inductance \cite{ammar2021led}. The series approximation is used for initial sizing; magnetizing current and burst dynamics are also considered when analyzing prototype switching and transients.
